% Phần 1: thay đổi từ v1 sang v2, năm nhãn, dữ liệu. Số liệu lấy qua \so{...} (so_lieu.tex).

\section{Từ v1 đến v2: thay đổi gì và vì sao}
\label{sec:p1-v1v2}

\textbf{v1 (vòng loại)} phân loại một câu bình luận thành HATE hoặc CLEAN. Câu đi qua bộ chuẩn hóa 5 tầng
(Unicode, ký tự đồng dạng, ký tự phân cách, ký tự lặp, teencode), rồi vào ViHateT5 đã fine-tune đối kháng (FGM).
Phần giải thích dùng occlusion: che lần lượt từng cụm hai từ, chạy lại mô hình, cụm nào làm điểm rủi ro giảm mạnh
thì tô sáng. Trên test của vòng loại (\so{v1goc.n} câu bình luận ViHSD/ViHOS), F1 lớp HATE của v1 là
\so{v1goc.C0} với câu sạch, \so{v1goc.C1} với câu né lọc (C1) và \so{v1goc.C2} với kiểu điểm mù (C2).

\textbf{Vấn đề.} Khi chạy lại v1 trên test mới của v2 (hội thoại do người viết, mục~\ref{sec:p1-test}), v1 báo
HATE cho \so{C0.v1.xp.r} số đoạn \nhan{XUC_PHAM}, nhưng chỉ \so{C0.v1.dd.r} đoạn đe dọa, \so{C0.v1.qr.r} đoạn
quấy rối và \so{C0.v1.ld.r} đoạn lừa đảo. Lý do nằm ở định nghĩa: HATE trong các bộ dữ liệu cũ gần với ``chửi bậy,
miệt thị''. Còn tin nhắn nguy hại thật thì khác: tin lừa đảo thường rất lịch sự, lời dọa có thể không có từ tục
nào. Vì vậy v2 thay đổi ba thứ:

\begin{itemize}
  \item \textbf{Nhãn theo tình huống.} 5 nhãn thay cho HATE/CLEAN, vì mỗi tình huống cần ứng dụng phản ứng khác
        nhau: làm mờ câu chửi, gợi ý chặn kẻ quấy rối, cảnh báo mạnh khi bị dọa, vô hiệu link lừa đảo.
  \item \textbf{Đọc cả đoạn hội thoại.} Ý đồ xấu thường chỉ lộ ở vài dòng cuối, sau mấy câu chào hỏi bình thường.
        Mô hình nhận tối đa 8 dòng gần nhất thay vì một câu.
  \item \textbf{Bằng chứng ra cùng lúc với nhãn.} Mô hình tự trả về các cụm từ làm bằng chứng trong cùng một lượt
        chạy, không phải chạy lại nhiều lần như occlusion. Nhãn và bằng chứng có ngưỡng riêng, được đo riêng.
\end{itemize}

\begin{table}[H]
\centering
\caption{So sánh v1 và v2}
\label{tab:p1-v1v2}
\small
\begin{tabularx}{\linewidth}{@{}p{2.6cm}YY@{}}
\toprule
 & \textbf{v1 (vòng loại)} & \textbf{v2 (chung kết)} \\
\midrule
Nhãn & HATE / CLEAN & 5 nhãn theo tình huống (mục~\ref{sec:p1-nhan}) \\
Đầu vào & Một câu bình luận & Đoạn hội thoại hai người, 1--8 dòng, mỗi dòng là một lượt nói \\
Mô hình & ViHateT5 (encoder--decoder), fine-tune đối kháng FGM & mmBERT-base, một bộ mã hóa với hai đầu ra: nhãn và
  cụm bằng chứng \\
Chống né lọc & Chuẩn hóa 5 tầng + FGM & Tăng cường dữ liệu khi train (bỏ dấu, biến dạng bề mặt) + FGM + nhánh đọc
  theo byte; chuẩn hóa thành tùy chọn và được đo bật/tắt \\
Giải thích & Occlusion sau khi phân loại, chạy lại mô hình nhiều lần & Gắn thẻ cụm bằng chứng ngay trong lượt
  chạy; cụm luôn chép nguyên văn từ tin nhắn \\
Dữ liệu train & ViHSD, ViHOS với nhãn gốc & Câu thật được LLM gán lại + hội thoại LLM sinh theo kịch bản
  (mục~\ref{sec:p1-dulieu}) \\
Dữ liệu test & Split test của ViHSD/ViHOS & Hội thoại do người trong nhóm tự viết, kèm bản biến dạng C1, C2 \\
Cách đo & F1 lớp HATE trên C0/C1/C2 & F1 từng nhãn; theo giọng vùng miền và kịch bản chưa thấy; precision/recall
  bằng chứng đo riêng; so với LLM lớn qua API \\
Dung lượng (FP32) & ONNX \so{v1.onnx} MiB & Trọng số \so{v2.trongso} MiB \\
\bottomrule
\end{tabularx}
\end{table}

\section{Năm nhãn mới}
\label{sec:p1-nhan}

Mỗi mẫu được gán đúng một nhãn trong bảng~\ref{tab:p1-nhan}. Định nghĩa đầy đủ và các ca khó nằm trong
\tep{docs/v2/HUONG_DAN_GAN_NHAN.md}; người gán, LLM gán sơ bộ và LLM baseline đều dùng chung tài liệu này.

\begin{table}[H]
\centering
\caption{Năm nhãn của v2}
\label{tab:p1-nhan}
\small
\begin{tabularx}{\linewidth}{@{}p{2.1cm}YYp{2.9cm}@{}}
\toprule
\textbf{Nhãn} & \textbf{Nghĩa} & \textbf{Ví dụ} & \textbf{Ứng dụng làm gì} \\
\midrule
\nhan{AN_TOAN} & Không thuộc 4 nhãn dưới & ``Đơn hàng của bạn đang được giao, vui lòng giữ máy'' & Hiện bình thường \\
\nhan{XUC_PHAM} & Lăng mạ, miệt thị một người hoặc một nhóm (kể cả vùng miền) & ``Đúng là não cá vàng, nói mãi
  không hiểu'' & Làm mờ, bấm mới xem \\
\nhan{QUAY_ROI} & Gạ gẫm không mong muốn, bám riết, bêu riếu đời tư & ``Sao seen mà không rep, anh nhắn 50 tin rồi
  đó'' & Gợi ý chặn, báo cáo \\
\nhan{DE_DOA} & Dọa gây hại thân thể, tài sản, danh dự; tống tiền & ``Mai tao chặn đầu mày ở cổng trường'' & Cảnh
  báo mạnh, gợi ý lưu bằng chứng \\
\nhan{LUA_DAO} & Thông tin giả để dụ chuyển tiền, đưa OTP, bấm link, cài ứng dụng & ``Chúc mừng bạn trúng iPhone,
  nộp phí ship 300k để nhận quà'' & Cảnh báo, vô hiệu link \\
\bottomrule
\end{tabularx}
\end{table}

\textbf{Khi một đoạn dính nhiều nhãn}, người gán hỏi lần lượt và dừng ở câu trả lời ``có'' đầu tiên:
\begin{enumerate}
  \item Có dọa gây hại, hoặc dọa để ép ai làm gì không? $\to$ \nhan{DE_DOA}
  \item Có dụ đưa tiền, OTP, mật khẩu, bấm link, cài ứng dụng bằng thông tin giả không? $\to$ \nhan{LUA_DAO}
  \item Có gạ gẫm, bình phẩm cơ thể, bám riết hoặc bêu riếu một người cụ thể không? $\to$ \nhan{QUAY_ROI}
  \item Có chửi, lăng mạ, miệt thị người hoặc nhóm người không? $\to$ \nhan{XUC_PHAM}
  \item Còn lại $\to$ \nhan{AN_TOAN}
\end{enumerate}
Ví dụ: ``Chuyển 10 triệu trước 8h không tao gửi clip cho cả nhà mày'' có đòi tiền, nhưng là tống tiền nên thuộc
\nhan{DE_DOA}, không phải \nhan{LUA_DAO}. Với đoạn nhiều dòng, nhãn của cả đoạn là tình huống nặng nhất xuất hiện
trong đoạn.

\textbf{Ca khó: trông nguy hại nhưng an toàn.} Đây là nguồn báo nhầm chính, nên được ghi rõ trong hướng dẫn và
được sinh thêm làm dữ liệu train (mục~\ref{sec:p1-sinh}):

\begin{table}[H]
\centering
\caption{Một số ca \nhan{AN_TOAN} dễ nhầm}
\label{tab:p1-cakho}
\small
\begin{tabularx}{\linewidth}{@{}p{6.4cm}Y@{}}
\toprule
\textbf{Tin nhắn} & \textbf{Vì sao an toàn} \\
\midrule
``Vãi cả chưởng, kèo này căng đét'' & Chửi thề cảm thán, không nhắm vào ai \\
``Tao giết mày luôn á, cười đau bụng'' & Nói quá kiểu đùa, dấu hiệu đùa nằm ngay trong câu \\
``Mã OTP của bạn là 123456. Không chia sẻ mã này cho bất kỳ ai'' & Tin OTP thật: không đòi gì, còn dặn đừng chia sẻ \\
``Ai bảo người miền Tây lười là phân biệt vùng miền đấy'' & Trích lại để phản bác \\
\bottomrule
\end{tabularx}
\end{table}

\textbf{Bằng chứng} là các cụm chép nguyên văn từ tin nhắn, cho biết vì sao đoạn đó nguy hại. Mỗi cụm nằm gọn
trong một dòng, lấy cụm ngắn nhất còn đủ nghĩa, được lấy nhiều cụm. \nhan{AN_TOAN} luôn rỗng; câu nguy hại cũng
được để rỗng nếu cái hại nằm ở cả câu (ví dụ mỉa mai). Ví dụ dưới đây là một đoạn thật trong test (giọng
\so{vd.dung.vung}, kịch bản \so{vd.dung.kichban}: dọa tìm ra ngoài đời trong game) và đầu ra của v2:

\ghinho{\textbf{Đầu vào} (mỗi dòng là một lượt nói, không ghi ai nói; nhãn đúng: \so{vd.dung.gold}):\\[2pt]
\so{vd.dung.dong}\\[4pt]
\textbf{Đầu ra của v2:}\\[2pt]
{\ttfamily \so{vd.dung.json}}}

\section{Dữ liệu}
\label{sec:p1-dulieu}

Dữ liệu v2 có ba phần, đi theo hai đường riêng (hình~\ref{fig:p1-luong}). Train gồm câu thật được LLM gán lại
theo 5 nhãn và hội thoại do LLM sinh theo kịch bản. Test do người trong nhóm tự viết, không có câu nào do LLM sinh.

\begin{figure}[H]
\centering
\begin{tikzpicture}[
  hop/.style={draw=navy, fill=nenxanh, rounded corners=2pt, align=center, font=\footnotesize,
              text width=2.55cm, minimum height=1.3cm, inner sep=3pt},
  hopdich/.style={hop, fill=navy!15, text width=1.9cm, font=\footnotesize\bfseries},
  hoptest/.style={hop, fill=nenxam, draw=chuxam},
  mui/.style={-{Stealth[length=2mm]}, thick, navy},
  node distance=5mm and 4.5mm]
% Nhánh 1: câu thật
\node[hop] (a1) {ViHSD, ViHOS, ViTHSD\\ \so{dl.goc.tong} bản ghi};
\node[hop, right=of a1] (a2) {Gộp, khử trùng\\ \so{dl.khacnhau} câu khác nhau};
\node[hop, right=of a2] (a3) {Tách kho test \so{dl.khotest} câu\\ (không gửi LLM)};
\node[hop, right=of a3] (a4) {LLM gán lại \so{dl.ganlai} câu; giữ \so{dl.that.giu} câu khớp người gán gốc};
% Nhánh 2: sinh theo kịch bản
\node[hop, below=of a1] (b1) {Kịch bản $\times$ kênh $\times$ giọng $\times$ độ lộ};
\node[hop, right=of b1] (b2) {LLM sinh \so{dl.sinh.tong} đoạn, kèm sinh đôi \nhan{AN_TOAN}};
\node[hop, right=of b2] (b3) {LLM gán lại, không biết nhãn yêu cầu};
\node[hop, right=of b3] (b4) {Giữ \so{dl.sinh.giu} đoạn trùng nhãn};
% Gộp thành train
\path (a4.east) -- (b4.east) coordinate[midway] (giua);
\node[hopdich, right=5mm of giua] (tr) {Train\\ \so{tr.tong} đoạn\\ {\normalfont(chọn cân nhãn)}};
\draw[mui] (a1) -- (a2); \draw[mui] (a2) -- (a3); \draw[mui] (a3) -- (a4);
\draw[mui] (b1) -- (b2); \draw[mui] (b2) -- (b3); \draw[mui] (b3) -- (b4);
\draw[mui] (a4.east) -- (tr.north west);
\draw[mui] (b4.east) -- (tr.south west);
% Nhánh 3: test do người viết
\node[hoptest, below=10mm of b1] (c1) {Người trong nhóm tự viết, trước khi đọc câu sinh};
\node[hoptest, right=of c1] (c2) {Test C0: \so{te.tong} đoạn, 3 giọng vùng};
\node[hoptest, right=of c2] (c3) {Sinh biến dạng C1 (\so{te.C1.n}) và C2 (\so{te.C2.n})};
\node[hoptest, right=of c3] (c4) {Kiểm tra rò rỉ: \so{rr.trunggan} cặp gần trùng với train};
\draw[-{Stealth[length=2mm]}, thick, chuxam] (c1) -- (c2);
\draw[-{Stealth[length=2mm]}, thick, chuxam] (c2) -- (c3);
\draw[-{Stealth[length=2mm]}, thick, chuxam] (c3) -- (c4);
\draw[dashed, chuxam] (tr.south) |- (c4.east);
\end{tikzpicture}
\caption{Luồng tạo dữ liệu v2. Hàng trên: câu thật; hàng giữa: hội thoại sinh; hàng dưới (xám): test do người viết.}
\label{fig:p1-luong}
\end{figure}

\subsection{Câu thật: gán lại bằng LLM}
\label{sec:p1-that}

\begin{itemize}
  \item \textbf{Gộp nguồn.} Ba bộ ViHSD (\so{dl.goc.vihsd}), ViHOS (\so{dl.goc.vihos}), ViTHSD (\so{dl.goc.vithsd})
        có tổng \so{dl.goc.tong} bản ghi. ViHOS và ViTHSD phần lớn dùng lại câu của ViHSD, nên sau khi gộp theo nội
        dung còn \so{dl.khacnhau} câu khác nhau.
  \item \textbf{Tách kho test trước khi gọi LLM.} Sau khử trùng gần đúng, các nhóm câu chưa từng nằm trong split train
        của ViHSD/ViHOS (phần ViHateT5 của v1 đã học) và chưa từng gửi LLM được tách thành kho test
        (\so{dl.khotest} câu). Kho này không bao giờ gửi cho LLM và không vào train.
  \item \textbf{Chọn câu cần gán lại.} Từ \so{dl.khotrain} câu còn lại, chọn \so{dl.ganlai} câu:
        \so{dl.ganlai.nguyhai} câu nhãn gốc nguy hại, \so{dl.ganlai.blacklist} câu sạch nhưng có từ trong blacklist
        của v1 (ca dễ báo nhầm), \so{dl.ganlai.tukhoa} câu khớp từ khóa của 3 nhãn hiếm, và \so{dl.ganlai.ngaunhien}
        câu sạch chọn ngẫu nhiên (seed 42).
  \item \textbf{Gán lại.} DeepSeek V4.1 Flash (qua OpenRouter, tắt chế độ suy nghĩ) đọc hướng dẫn gán nhãn và trả về
        nhãn kèm bằng chứng dạng JSON. Script tự kiểm nhãn hợp lệ và bỏ cụm bằng chứng không có nguyên văn trong câu.
  \item \textbf{Lọc theo đồng thuận.} Chỉ giữ câu mà nhãn LLM (xét sạch/nguy hại) khớp với ít nhất một người gán
        gốc. Giữ \so{dl.that.giu} câu, bỏ \so{dl.that.bo} câu. Phần lớn câu bị bỏ là \so{dl.that.bo.hate.antoan} câu
        gốc HATE mà LLM cho là an toàn và \so{dl.that.bo.clean.xucpham} câu gốc CLEAN mà LLM cho là xúc phạm. Nhóm
        không đủ người để xét tay từng câu mâu thuẫn, nên chọn bỏ chúng thay vì giữ một nhãn có thể sai.
\end{itemize}

\subsection{Hội thoại sinh theo kịch bản}
\label{sec:p1-sinh}

Câu thật gần như không có ba nhãn hiếm: sau khi gán lại chỉ còn \so{tr.dd.that} câu \nhan{DE_DOA},
\so{tr.qr.that} câu \nhan{QUAY_ROI} và \so{tr.ld.that} câu \nhan{LUA_DAO}. Vì vậy nhóm cho LLM sinh hội thoại theo
kịch bản (\tep{docs/v2/KICH_BAN.md}).

\begin{itemize}
  \item \textbf{Kịch bản.} 26 kịch bản xoay quanh học sinh, sinh viên: 12 lừa đảo (nạp game giá rẻ, việc làm thêm
        phải nộp cọc, giveaway giả...), 6 đe dọa (hẹn đánh, đòi nợ, tống tiền bằng ảnh...), 6 quấy rối (xin ảnh
        riêng tư, nhắn dồn dập sau khi bị từ chối...), 2 xúc phạm. Bốn kịch bản (\so{rr.kichban}) được giữ ngoài
        train, dùng làm phần ``kịch bản chưa thấy'' khi test.
  \item \textbf{Ô sinh.} Mỗi lượt gọi LLM ứng với một ô = kịch bản $\times$ kênh $\times$ giọng $\times$ độ lộ
        $\times$ kiểu lộ, và xin 5 đoạn khác nhau. Kênh: Messenger/Zalo, nhóm lớp, chat game (2--6 dòng); bình
        luận, SMS (1 dòng). Giọng: teencode, không dấu, viết tắt, giọng Bắc, Trung, Nam. Độ lộ: lộ liễu hoặc ẩn
        ý. Kiểu lộ: lộ ngay dòng đầu hoặc lộ dần ở cuối.
  \item \textbf{Sinh đôi an toàn.} Mỗi ô nguy hại có một ô \nhan{AN_TOAN} cùng kênh, cùng giọng: đoạn trông giống
        nguy hại nhưng vô hại (bạn thân đùa, nhắc trả tiền trà sữa lịch sự). Mục đích là để mô hình không học mẹo
        ``cứ có chữ dọa là nguy hại''. \nhan{XUC_PHAM} chỉ sinh dạng nhiều dòng, vì câu \nhan{XUC_PHAM} thật đều một
        dòng; như vậy mô hình không học được mẹo ``nhiều dòng là nhãn hiếm''.
  \item \textbf{An toàn nội dung.} Không có nội dung tình dục tường minh, không tên, số điện thoại, tài khoản thật;
        link là tên miền giả dạng \tep{xxx[.]com}. Ô nào LLM từ chối thì bỏ qua.
  \item \textbf{Kết quả sinh.} \so{cp.sinh.goi} lượt gọi cho \so{dl.sinh.tong} đoạn, trung bình \so{dl.sinh.dongtb}
        dòng mỗi đoạn.
  \item \textbf{Lọc lại.} LLM gán nhãn lại từng đoạn mà không biết nhãn đã yêu cầu khi sinh; chỉ giữ đoạn trùng
        nhãn: \so{dl.sinh.giu} đoạn (\so{dl.sinh.tilegiu}). \nhan{QUAY_ROI} giữ ít nhất (\so{dl.sinh.qr.tilegiu}),
        chủ yếu do kịch bản bình phẩm ngoại hình (QR3) chỉ giữ \so{dl.sinh.qrba.tilegiu}: các đoạn này hay bị gán lại
        thành \nhan{XUC_PHAM}. Ranh giới giữa hai nhãn này là chỗ hướng dẫn gán nhãn còn chưa rõ.
  \item Đoạn sinh được khử trùng gần đúng, và đoạn nào gần trùng câu test người viết thì bị loại.
\end{itemize}

\subsection{Tập train và tập test}
\label{sec:p1-test}

\begin{figure}[H]
\centering
% Sinh tự động bằng `python -m src.v2.eval.xuat_so_lieu_bao_cao`. Không sửa tay.
\begin{tikzpicture}
\begin{axis}[xbar stacked, bar width=11pt, width=0.9\linewidth, height=5.2cm,
  symbolic y coords={ld,dd,qr,xp,at}, ytick=data, yticklabels={\nhan{LUA_DAO},\nhan{DE_DOA},\nhan{QUAY_ROI},\nhan{XUC_PHAM},\nhan{AN_TOAN}},
  xmin=0, xlabel={Số mẫu train}, axis x line*=bottom, axis y line*=left, xmajorgrids,
  grid style={draw=black!12}, tick label style={font=\small}, label style={font=\small},
  scaled x ticks=false, xticklabel style={/pgf/number format/.cd, 1000 sep={.}},
  legend style={at={(0.5,1.03)}, anchor=south, legend columns=2, draw=none, font=\small,
    /tikz/every even column/.append style={column sep=0.4cm}}]
\addplot[fill=navy, draw=white, area legend] coordinates {(0,ld) (156,dd) (62,qr) (2527,xp) (3967,at)};
\addlegendentry{Câu thật (gán lại)}
\addplot[fill=maullm, draw=white, area legend] coordinates {(1751,ld) (1361,dd) (1098,qr) (470,xp) (2000,at)};
\addlegendentry{Đoạn LLM sinh}
\end{axis}
\end{tikzpicture}

\caption{Thành phần tập train theo nhãn. Ba nhãn hiếm (\nhan{QUAY_ROI}, \nhan{DE_DOA}, \nhan{LUA_DAO}) gần như toàn
là đoạn LLM sinh, nên điểm của chúng phải đo lại trên câu người viết.}
\label{fig:p1-train}
\end{figure}

\begin{table}[H]
\centering
\caption{Số mẫu theo nhãn. Train lưu bản gốc C0; test C0 do người viết}
\label{tab:p1-train}
\small
\begin{tabular}{@{}lrrrr@{}}
\toprule
\textbf{Nhãn} & \textbf{Train: câu thật} & \textbf{Train: đoạn sinh} & \textbf{Train: tổng} & \textbf{Test C0} \\
\midrule
\nhan{AN_TOAN}  & \so{tr.at.that} & \so{tr.at.sinh} & \so{tr.at} & \so{te.at} \\
\nhan{XUC_PHAM} & \so{tr.xp.that} & \so{tr.xp.sinh} & \so{tr.xp} & \so{te.xp} \\
\nhan{QUAY_ROI} & \so{tr.qr.that} & \so{tr.qr.sinh} & \so{tr.qr} & \so{te.qr} \\
\nhan{DE_DOA}   & \so{tr.dd.that} & \so{tr.dd.sinh} & \so{tr.dd} & \so{te.dd} \\
\nhan{LUA_DAO}  & \so{tr.ld.that} & \so{tr.ld.sinh} & \so{tr.ld} & \so{te.ld} \\
\midrule
Tổng & \so{tr.that} & \so{tr.sinh} & \so{tr.tong} & \so{te.tong} \\
\bottomrule
\end{tabular}
\end{table}

\textbf{Train} có \so{tr.tong} đoạn. Ba nhãn hiếm gần như toàn là đoạn sinh, nên điểm thật của chúng chỉ đo được
trên test. Câu thật đều một dòng (\so{tr.motdong} mẫu một dòng), còn đoạn sinh phần lớn nhiều dòng; phần kết quả
kiểm tra mô hình không đoán nhãn theo số dòng.

\textbf{Test} có \so{te.tong} đoạn do \so{te.nguoiviet} người trong nhóm tự viết theo
\tep{docs/v2/HUONG_DAN_VIET_TEST.md}, viết xong trước khi đọc bất kỳ đoạn LLM sinh nào để không vô tình bắt chước.
\begin{itemize}
  \item \textbf{Giọng vùng miền:} Bắc \so{te.bac}, Trung \so{te.trung}, Nam \so{te.nam} đoạn, viết đúng cách người
        mỗi miền nhắn (``tau nỏ cần mô'', ``tui hổng cần đâu'').
  \item \textbf{Kịch bản:} \so{te.trongpp} đoạn thuộc kịch bản đã sinh vào train; \so{te.chuathay} đoạn thuộc 4 kịch
        bản giữ ngoài train và các tình huống mới do người viết tự nghĩ. Phần chưa thấy chỉ có \so{te.chuathay.xp}
        đoạn \nhan{XUC_PHAM}.
  \item \textbf{Nhãn:} người viết gán nhãn và bằng chứng; một LLM gán độc lập để đối chiếu. Hai bên khớp ở
        \so{te.chon.trung} đoạn. Ở các đoạn lệch, người duyệt giữ nhãn người viết cho \so{te.chon.nguoi} đoạn và đổi
        theo LLM cho \so{te.chon.llm} đoạn (ghi \tep{cach_gan=llm_nguoi_duyet}). Test chưa có lượt gán chéo giữa hai
        người nên chưa có kappa.
\end{itemize}

\textbf{Bản biến dạng C1, C2.} Mỗi đoạn test được biến dạng thành 2 bản cho mỗi nhóm, kiểu bốc ngẫu nhiên (seed 42),
được \so{te.C1.n} mẫu C1 và \so{te.C2.n} mẫu C2. Biến dạng làm theo từng từ, khoảng \so{te.C1.tiletu} số từ, luôn
trúng cụm bằng chứng, và vị trí bằng chứng đi theo từ bị đổi.

\begin{table}[H]
\centering
\caption{Các kiểu biến dạng trong test (số mẫu)}
\label{tab:p1-bien-dang}
\small
\begin{tabularx}{\linewidth}{@{}p{3.6cm}Y@{}}
\toprule
\textbf{Nhóm} & \textbf{Kiểu} \\
\midrule
C1, đã thấy (xuất hiện tự nhiên trong train) & bỏ dấu (\so{te.C1.bodau}), teencode (\so{te.C1.teencode}),
  gõ sai (\so{te.C1.gosai}), kéo dài ký tự (\so{te.C1.keodai}) \\
C1, giữ ngoài train & dính từ (\so{te.C1.dinhtu}), số thay chữ (\so{te.C1.sothaychu}),
  chèn ký tự (\so{te.C1.chenkytu}) \\
C2, điểm mù, không có trong train & chữ toàn độ rộng (\so{te.C2.fullwidth}), ký tự vô hình (\so{te.C2.zerowidth}),
  chèn dấu cách giữa các chữ (\so{te.C2.spaceinjection}), ký tự phân cách lạ (\so{te.C2.exoticseparator}),
  chữ Cyrillic giống chữ Latin (\so{te.C2.cyrillichomoglyph}), kéo dài kép (\so{te.C2.doubleelongation}) \\
\bottomrule
\end{tabularx}
\end{table}

\textbf{Chống rò rỉ.} Script \tep{src.v2.data.kiem_tra_ro_ri} so test với train và kho test: \so{rr.trunggan} cặp gần
trùng (điểm trùng 3-gram từ \so{rr.nguong} trở lên, so cả đoạn và từng dòng), \so{rr.trainkichban} câu train thuộc
4 kịch bản giữ ngoài, không trùng id. Các kiểu C2 và nhóm C1 giữ ngoài không được dùng khi tăng cường dữ liệu lúc
train.

\subsection{Chất lượng nhãn train: điều chưa làm được}
\label{sec:p1-chatluong}

Train chủ yếu do LLM gán hoặc sinh, nên cần ước tỉ lệ nhãn sai. Kế hoạch là hai người gán độc lập mẫu
\so{dt.mau} câu (một nửa câu được giữ, một nửa câu bị bỏ, che nhãn). \textbf{Thực tế mẫu này do một mô hình thứ
hai (Claude) gán, chưa phải người.} Vì vậy chưa có kappa giữa người, và các số sau chỉ là mức đồng ý giữa hai mô hình:

\begin{itemize}
  \item Ở nhóm câu được giữ, mô hình thứ hai đồng ý với nhãn của DeepSeek ở \so{dt.giu.dung}/\so{dt.giu.tong} câu
        (\so{dt.giu.tile}, khoảng tin cậy 95\%: \so{dt.giu.lo} -- \so{dt.giu.hi}). Như vậy vẫn còn một phần nhãn train
        có thể sai.
  \item Ở nhóm câu bị bỏ, mô hình thứ hai đứng về phía LLM (thay vì nhãn gốc) trong \so{dt.bo.llm} trường hợp: bỏ
        câu mâu thuẫn cũng làm mất một số câu mà LLM gán đúng.
\end{itemize}
Do đó điểm trên test người viết là thước đo chính, và nhãn train nhiễu được ghi là một hạn chế.

\subsection{Chi phí tạo dữ liệu}
\label{sec:p1-chiphi}

\begin{table}[H]
\centering
\caption{Chi phí gọi LLM để tạo dữ liệu train (DeepSeek V4.1 Flash, tính từ số token $\times$ giá endpoint)}
\label{tab:p1-chiphi}
\small
\begin{tabular}{@{}lrrr@{}}
\toprule
\textbf{Bước} & \textbf{Lượt gọi} & \textbf{USD} & \textbf{Phút} \\
\midrule
Gán lại câu thật & \so{cp.ganlai.goi} & \so{cp.ganlai.usd} & \so{cp.ganlai.phut} \\
Sinh hội thoại & \so{cp.sinh.goi} & \so{cp.sinh.usd} & \so{cp.sinh.phut} \\
Lọc lại đoạn sinh & \so{cp.loc.goi} & \so{cp.loc.usd} & \so{cp.loc.phut} \\
\midrule
Tổng & \so{cp.tong.goi} & \so{cp.tong.usd} & \so{cp.tong.phut} \\
\bottomrule
\end{tabular}
\end{table}

Toàn bộ phần dữ liệu do LLM tạo và gán tốn \so{cp.tong.usd} USD (khoảng \so{cp.tong.vnd} đồng theo tỷ giá giả định
ghi trong \tep{chi_phi.json}), khoảng \so{cp.tong.phut} phút chạy và \so{cp.tong.trieutoken} triệu token. Chi phí
thấp nên có thể sinh lại hoặc mở rộng kịch bản khi xuất hiện kiểu lừa đảo mới.
